% !TEX root = ../main.tex
\corrigesChapitre{27}{Mesures et incertitudes}
\corrigesSouples

\begin{blocLivre}{cor-27-1}
\begin{corrige}{1}{Moyenne et écart-type d'une série}
\demarche{Moyenne, écart-type expérimental (avec $N-1$), puis incertitude-type de la moyenne $u=s/\sqrt N$.}
\begin{questions}
\item Le lâcher du pendule, le déclenchement du chronomètre et les conditions de l'expérience varient légèrement d'une
mesure à l'autre : c'est la \res{\text{variabilité}} de la mesure.
\item $\bar T=\dfrac{11{,}48}{8}=\res{\qty{1,435}{s}}$.
\item La somme des carrés des écarts à la moyenne vaut \qty{1,8e-3}{s^2} : $s=\sqrt{\dfrac{\num{1.8e-3}}{7}}=\res{\qty{0,016}{s}}$.
\item $u(\bar T)=\dfrac{0{,}016}{\sqrt8}=\qty{0,0057}{s}$ : $\res{T=\qty{1,4350}{s}}$ avec $\res{u(T)=\qty{0,0057}{s}}$.
\end{questions}
\end{corrige}
\end{blocLivre}
\begin{blocLivre}{cor-27-2}
\begin{corrige}{2}{Une seule mesure : évaluation de type B}
\demarche{Identifier la demi-étendue $a$, puis $u=a/\sqrt3$ ; incertitude relative $u/m$.}
\begin{questions}
\item $a=\qty{0,5}{\degreeCelsius}$ : $u=\dfrac{0{,}5}{\sqrt3}=\res{\qty{0,29}{\degreeCelsius}}$.
\item $u=\dfrac{0{,}15}{\sqrt3}=\res{\qty{0,087}{mL}}$.
\item $a=0{,}008\times4{,}82+2\times0{,}01=\res{\qty{0,059}{V}}$, d'où $u=\dfrac{0{,}059}{\sqrt3}=\res{\qty{0,034}{V}}$.
\item $\dfrac{0{,}29}{21}=1{,}4~\%$ ; $\dfrac{0{,}087}{250}=0{,}035~\%$ ; $\dfrac{0{,}034}{4{,}82}=0{,}71~\%$. La \res{\text{fiole jaugée}} donne la mesure la
plus précise.
\end{questions}
\end{corrige}
\end{blocLivre}
\begin{blocLivre}{cor-27-3}
\begin{corrige}{3}{Écrire un résultat}
\demarche{Arrondir d'abord l'incertitude-type, puis arrêter la valeur au même rang que son dernier chiffre.}
\begin{questions}
\item On arrondit chaque incertitude-type à deux chiffres significatifs, puis la valeur au même rang : voir la colonne
« Résultat » du tableau.
\item Voir la dernière colonne du tableau.
\item L'incertitude porte sur les dixièmes : les chiffres suivants de la valeur n'ont pas de sens. On écrit
$\res{g=\qty{9,8}{m.s^{-2}}}$ avec $u(g)=\qty{0,1}{m.s^{-2}}$.
\end{questions}
\begin{center}\small
\begin{tabular}{lccc}
\toprule
Grandeur & Incertitude-type & Résultat & Relative \\
\midrule
masse & \qty{0,024}{g} & \res{\qty{12,348}{g}} & 0,19~\% \\
concentration & \qty{0,00087}{mol.L^{-1}} & \res{\qty{0,05124}{mol.L^{-1}}} & 1,7~\% \\
longueur d'onde & \qty{1,5}{nm} & \res{\qty{632,8}{nm}} & 0,24~\% \\
célérité & \qty{3,4}{m.s^{-1}} & \res{\qty{343,3}{m.s^{-1}}} & 0,99~\% \\
\bottomrule
\end{tabular}
\end{center}
\end{corrige}
\end{blocLivre}
\begin{blocLivre}{cor-27-4}
\begin{corrige}{4}{Lire un histogramme}
\demarche{Lire les effectifs ; moyenne d'une série regroupée en classes : centres des classes pondérés par les effectifs.}
\begin{questions}
\item La classe \res{[11{,}9~;~12{,}0[} (15 binômes). $15+14=29$ binômes sur 50, soit \res{\qty{58}{\percent}}.
\item La somme des produits « centre $\times$ effectif » vaut $11{,}75\times2+11{,}85\times7+\dots+12{,}25\times3=\qty{588,35}{mL}$,
d'où $\bar V\approx\dfrac{588{,}35}{49}=\res{\qty{12,01}{mL}}$.
\item Une \res{\text{équivalence dépassée}} (versement trop rapide, changement de couleur mal repéré) ou une erreur de lecture de
la burette (\qty{13,4}{mL} au lieu de \qty{12,4}{mL}).
\item $\bar V\approx\dfrac{588{,}35+13{,}45}{50}=\res{\qty{12,04}{mL}}$. On ne l'écarte que si l'erreur est \res{\text{identifiée}} (cahier de
laboratoire, nouvelle mesure) ; sinon, on la conserve en la signalant.
\end{questions}
\end{corrige}
\end{blocLivre}
\begin{blocLivre}{cor-27-5}
\begin{corrige}{5}{Comparer à une valeur de référence}
\demarche{$z=\dfrac{|c_\mathrm{mes}-c_\mathrm{ref}|}{u(c)}$ ; compatible si $z<2$.}
\begin{questions}
\item $z_A=\dfrac{0{,}16}{0{,}12}=\res{1{,}3}$ ; $z_B=\dfrac{0{,}33}{0{,}08}=\res{4{,}1}$ ; $z_C=\dfrac{0{,}07}{0{,}02}=\res{3{,}5}$.
\item Seule la mesure du \res{\text{groupe A}} est compatible avec la référence.
\item Son incertitude-type est très faible : l'écart de $0{,}07$ représente 3,5 incertitudes-types. Une source d'incertitude
a sans doute été \res{\text{oubliée}}, ou une erreur systématique décale la mesure : des pertes thermiques réduisent
l'élévation de température et font surestimer $c$.
\end{questions}
\end{corrige}
\end{blocLivre}
\begin{blocLivre}{cor-27-6}
\begin{corrige}{6}{Chronométrer un pendule à la main}
\demarche{Une mesure : $u=s$ ; moyenne de $N$ mesures : $u=s/\sqrt N$ ; une durée de 20 périodes divisée par 20.}
\begin{questions}
\item $u(T)=\qty{0,12}{s}$, soit $\dfrac{0{,}12}{1{,}43}=\res{\qty{8,4}{\percent}}$.
\item $u(\bar T)=\dfrac{0{,}12}{\sqrt{10}}=\res{\qty{0,038}{s}}$.
\item $\sqrt N=\dfrac{0{,}12}{0{,}010}=12$, soit $\res{N=144}$ mesures : la méthode est bien trop longue.
\item L'erreur de déclenchement ne porte qu'une fois sur la durée totale : $u(T)=\dfrac{0{,}12}{20}=\res{\qty{0,006}{s}}$, et
$\dfrac{0{,}038}{20}=\res{\qty{0,0019}{s}}$ avec dix mesures. Mesurer une durée longue est bien plus efficace que multiplier les
mesures.
\end{questions}
\end{corrige}
\end{blocLivre}
\begin{blocLivre}{cor-27-7}
\begin{corrige}{7}{Cuivre ou laiton ?}
\demarche{$\rho=m/V$ ; composition des incertitudes relatives ; écarts normalisés avec chaque métal.}
\begin{questions}
\item $\rho=\dfrac{85{,}32}{9{,}5}=\res{\qty{8,98}{g.cm^{-3}}}$ (\qty{1}{mL} $=$ \qty{1}{cm^3}).
\item $\dfrac{u(m)}{m}=0{,}012~\%$ ; $\dfrac{u(V)}{V}=3{,}2~\%$ : le \res{\text{volume}} domine largement.
\item $u(\rho)=0{,}032\times8{,}98=\res{\qty{0,28}{g.cm^{-3}}}$. $z_\text{cuivre}=\dfrac{0{,}02}{0{,}28}=0{,}07$ et $z_\text{laiton}=\dfrac{0{,}48}{0{,}28}=1{,}7$ : les deux
métaux sont compatibles, on \res{\text{ne peut pas conclure}}.
\item Il faudrait $z_\text{laiton}\geqslant2$, soit $u(\rho)<\dfrac{0{,}48}{2}=\res{\qty{0,24}{g.cm^{-3}}}$, donc
$u(V)<\dfrac{0{,}24}{8{,}98}\times9{,}5=\res{\qty{0,25}{mL}}$ : une éprouvette plus fine ou une pesée dans l'eau.
\end{questions}
\end{corrige}
\end{blocLivre}
\begin{blocLivre}{cor-27-8}
\begin{corrige}{8}{Le résultat d'un titrage}
\demarche{Type B pour la pipette ; $c_A=c_B\,V_\mathrm{E}/V_A$ ; composition des trois incertitudes relatives.}
\begin{questions}
\item $u(V_A)=\dfrac{0{,}02}{\sqrt3}=\res{\qty{0,012}{mL}}$.
\item $c_A=\dfrac{0{,}100\times12{,}35}{10{,}00}=\res{\qty{0,1235}{mol.L^{-1}}}$.
\item Les trois termes valent 1,0~\%, 0,32~\% et 0,12~\% : $\dfrac{u(c_A)}{c_A}=1{,}06~\%$ et $\res{u(c_A)=\qty{0,0013}{mol.L^{-1}}}$.
La concentration de la \res{\text{solution titrante}} domine.
\item $z=\dfrac{0{,}0015}{0{,}0013}=1{,}2<2$ : le résultat est \res{\text{compatible}} avec l'étiquette.
\end{questions}
\end{corrige}
\end{blocLivre}
\begin{blocLivre}{cor-27-9}
\begin{corrige}{9}{Un histogramme en Python}
\demarche{\code{mean} : moyenne ; \code{stdev} : écart-type expérimental ; \code{s / N**0.5} : incertitude-type de type A.}
\begin{questions}
\item \code{mean} calcule la moyenne et \code{stdev} l'écart-type expérimental (avec $N-1$). La dernière valeur est
$s/\sqrt N$, l'\res{\text{incertitude-type de la moyenne}}.
\item $\res{f'=\qty{100,5}{mm}}$ avec $\res{u(f')=\qty{0,60}{mm}}$ ($2{,}07/\sqrt{12}=0{,}598$).
\item Les bornes vont de 96 à 105 : \res{\text{9 classes}} de \qty{1}{mm}, de $[96~;~97[$ à $[104~;~105]$.
\item $z=\dfrac{0{,}5}{0{,}60}=0{,}83<2$ : la mesure est \res{\text{compatible}} avec la valeur du fabricant.
\end{questions}
\end{corrige}
\end{blocLivre}
\begin{blocLivre}{cor-27-10}
\begin{corrige}{10}{Une longueur d'onde par les fentes d'Young}
\demarche{$i=\dfrac{10\,i}{10}$ ; type B pour $a$ ; composition des trois incertitudes relatives ; écart normalisé.}
\begin{questions}
\item L'incertitude de repérage des franges est divisée par 10 : $\res{i=\qty{4,75}{mm}}$ et $\res{u(i)=\qty{0,05}{mm}}$.
\item $u(a)=\dfrac{0{,}010}{\sqrt3}=\res{\qty{0,0058}{mm}}$ ; $\lambda=\dfrac{\num{0.200e-3}\times\num{4.75e-3}}{1{,}500}=\res{\qty{633}{nm}}$.
\item Les termes valent 2,9~\%, 1,1~\% et 0,33~\% : $\dfrac{u(\lambda)}{\lambda}=3{,}1~\%$ et $\res{u(\lambda)=\qty{20}{nm}}$. L'\res{\text{écartement des fentes}}
domine.
\item $z=\dfrac{|633-632{,}8|}{20}=0{,}01<2$ : la mesure est \res{\text{compatible}} avec la longueur d'onde du laser.
\end{questions}
\end{corrige}
\end{blocLivre}
\begin{blocLivre}{cor-27-11}
\begin{corrige}{11}{Simuler une mesure de résistance}
\demarche{Une notice donne un intervalle sans valeur privilégiée : tirage uniforme ; l'écart-type des valeurs simulées est $u(R)$.}
\begin{questions}
\item Les notices donnent seulement un intervalle, dans lequel aucune valeur n'est privilégiée : on tire les valeurs
\res{\text{uniformément}} dans cet intervalle.
\item \code{R} contient \num{100000} valeurs simulées de $U/I$. La moyenne, \res{\qty{227,6}{\ohm}}, est la valeur retenue ; l'écart-type,
\res{\qty{2,1}{\ohm}}, est l'incertitude-type composée.
\item $u(U)=\dfrac{0{,}045}{\sqrt3}=\qty{0,026}{V}$ et $u(I)=\dfrac{0{,}3}{\sqrt3}=\qty{0,17}{mA}$ ;
$\dfrac{u(R)}{R}=\sqrt{\left(\dfrac{0{,}026}{5{,}03}\right)^2+\left(\dfrac{0{,}17}{22{,}1}\right)^2}=0{,}93~\%$, d'où $\res{u(R)=\qty{2,1}{\ohm}}$ : même résultat.
\item La tolérance définit l'intervalle $[\qty{209}{\ohm}~;~\qty{231}{\ohm}]$, qui contient \qty{227,6}{\ohm} : la résistance est \res{\text{conforme}}.
\end{questions}
\end{corrige}
\end{blocLivre}
\begin{blocLivre}{cor-27-12}
\begin{corrige}{12}{Étalonner une pipette jaugée}
\demarche{Type A sur les masses ; $V=m/\rho$ ; comparer à la tolérance de la classe A et calculer l'écart normalisé.}
\begin{questions}
\item $\bar m=\dfrac{99{,}805}{10}=\res{\qty{9,9805}{g}}$ et $\res{s=\qty{0,0040}{g}}$.
\item $u(\bar m)=\dfrac{0{,}0040}{\sqrt{10}}=\res{\qty{0,0013}{g}}$ : évaluation de \res{\text{type A}}, par une approche statistique.
\item $V=\dfrac{9{,}9805}{0{,}9982}=\res{\qty{9,9985}{mL}}$ et $\res{u(V)=\qty{0,0013}{mL}}$.
\item Le volume est compris entre \qty{9,98}{mL} et \qty{10,02}{mL} : la pipette est \res{\text{conforme}} à la classe A. De plus,
$z=\dfrac{0{,}0015}{0{,}0013}=1{,}2<2$ : le volume délivré est \res{\text{compatible}} avec \qty{10,00}{mL}.
\end{questions}
\end{corrige}
\end{blocLivre}
\begin{blocLivre}{cor-27-13}
\begin{corrige}{13}{Le comprimé de vitamine C}
\demarche{À l'équivalence, $n(\text{acide ascorbique})=n(\ce{I2})$ ; facteur de dilution 20 ; composition de quatre incertitudes relatives.}
\begin{questions}
\item À l'équivalence, les réactifs ont été introduits dans les proportions stœchiométriques (1 pour 1) :
$\res{n_\mathrm{titrée}=c\,V_\mathrm{E}}$.
\item La fiole contient 20 fois la quantité titrée : $m=20\,c\,V_\mathrm{E}\,M=20\times\num{5.00e-3}\times\num{28.0e-3}\times176{,}1=\res{\qty{493,1}{mg}}$.
\item $\dfrac{u(m)}{m}=\sqrt{1{,}0^2+0{,}18^2+0{,}12^2+0{,}044^2}~\%=1{,}02~\%$ : $\res{u(m)=\qty{5,0}{mg}}$. Il faut réduire en priorité
l'incertitude sur la \res{\text{concentration du diiode}}.
\item $z=\dfrac{500-493{,}1}{5{,}0}=1{,}4<2$ : la masse mesurée est compatible avec \qty{500}{mg}, et elle appartient à
l'intervalle toléré $[\qty{475}{mg}~;~\qty{525}{mg}]$ : le comprimé est \res{\text{conforme}}.
\end{questions}
\end{corrige}
\end{blocLivre}
\begin{blocLivre}{cor-27-14}
\begin{corrige}{14}{Mesurer g avec un pendule}
\demarche{$g=\dfrac{4\pi^2L}{T^2}$ ; type A sur la durée de 20 périodes, divisée par 20 ; composition fournie.}
\begin{questions}
\item $T^2=4\pi^2\,\dfrac{L}{g}$, d'où $\res{g=\dfrac{4\pi^2L}{T^2}}$.
\item $\overline{20\,T}=\qty{28,640}{s}$, $s=\qty{0,051}{s}$ et $u(\overline{20\,T})=\dfrac{0{,}051}{\sqrt{10}}=\qty{0,016}{s}$ : $\res{T=\qty{1,4320}{s}}$ avec
$\res{u(T)=\qty{0,00080}{s}}$.
\item $g=\dfrac{4\pi^2\times0{,}510}{1{,}4320^2}=\res{\qty{9,818}{m.s^{-2}}}$.
\item Les termes valent $\dfrac{0{,}003}{0{,}510}=0{,}59~\%$ et $2\times\dfrac{0{,}00080}{1{,}4320}=0{,}11~\%$ : $\res{u(g)=\qty{0,059}{m.s^{-2}}}$ ;
$z=\dfrac{0{,}008}{0{,}059}=0{,}14<2$ : mesure \res{\text{compatible}}. Il faut améliorer la mesure de la \res{\text{longueur}}, jusqu'au centre de
la masse.
\end{questions}
\end{corrige}
\end{blocLivre}
\begin{blocLivre}{cor-27-15}
\begin{corrige}{15}{Verre crown ou verre flint ?}
\demarche{$n=\dfrac{\sin i}{\sin r}$ ; la simulation tire chaque angle uniformément ; écarts normalisés avec les deux verres.}
\begin{questions}
\item $n=\dfrac{\sin\ang{40}}{\sin\ang{25}}=\dfrac{0{,}6428}{0{,}4226}=\res{1{,}521}$.
\item Chaque angle est connu à $\pm\ang{1}$ près, sans valeur privilégiée dans cet intervalle : on le tire
\res{\text{uniformément}} entre \ang{39} et \ang{41} (et entre \ang{24} et \ang{26}), puis on calcule l'indice.
\item $\res{n=1{,}522}$ avec $\res{u(n)=0{,}038}$. L'histogramme, aplati en plateau, n'a pas la forme en cloche habituelle :
il vient de tirages uniformes.
\item $z_\text{crown}=\dfrac{0{,}002}{0{,}038}=0{,}05<2$ et $z_\text{flint}=\dfrac{0{,}098}{0{,}038}=2{,}6>2$ : le verre est un \res{\text{crown}}.
\item Affiner le faisceau, ou mesurer plusieurs couples $(i,r)$ et tracer $\sin i$ en fonction de $\sin r$ : le coefficient
directeur de la droite donne $n$ avec une incertitude plus faible.
\end{questions}
\schema[corrige]{27-ex15}
\end{corrige}
\end{blocLivre}
